\section{Introduction}

Parsers are among the most exposed pieces of code in any system: they are
the first code to touch untrusted input, and a memory-safety bug in a parser
is directly attacker-controlled. Fuzzing -- automatically generating inputs
and observing how a target handles them -- is the standard defense, but the
two dominant families of fuzzer make different tradeoffs around structure and
feedback. Coverage-guided, mutation-based fuzzers such as AFL discover
interesting inputs by instrumenting the target and rewarding inputs that
reach new code, but for a structured format such as JSON, random byte-level
mutation of a seed corpus destroys syntactic validity far more often than it
discovers it, so most of the fuzzer's budget is spent on inputs the parser
rejects in its first few bytes. Grammar-based and property-based approaches
invert this: a generator is written directly against a grammar or schema, so
every input is syntactically valid (or deliberately near-valid) by
construction, and the fuzzer's budget goes toward exploring the target's
\emph{structural} handling of the format rather than re-discovering its
lexer. Hypothesis, the property-based testing library we build on, is one
such approach: a developer composes strategies -- generators for values of a
given shape -- that can mirror a grammar directly.

The gap is that a hand-written grammar-based generator is usually static. It
is authored once, and unless a human revisits it, it keeps producing the
same distribution of shapes for the rest of a campaign, no matter what the
target's actual behavior turns out to be. Coverage-guided fuzzers close an
analogous loop by using instrumentation feedback to reward and retain
interesting inputs; grammar-based generators typically have nothing
equivalent, because the whole appeal of generating from a grammar is that it
does not require instrumenting or even recompiling the target. This is not
a minor convenience: many realistic targets -- vendored third-party
libraries pinned at a specific commit, closed-source or license-restricted
binaries, embedded or cross-compiled parsers -- cannot be rebuilt with
coverage instrumentation at all, or the instrumented build would not be
representative of what actually ships. A feedback loop for grammar-based
fuzzing is only useful in practice if it can run under exactly this
constraint: black-box, with no coverage signal available.

This paper asks a narrow, empirical question under that constraint: given
only a formal grammar and the kind of feedback a black-box harness can
actually produce -- whether an input was accepted or rejected, what the
parser's own rejection reason was, and whether a sanitizer fired -- can a
large language model (LLM) act as an effective iterative refiner of a
grammar-based generator, without ever seeing which lines of the target ran?
We answer it by building a complete pipeline around the cJSON parser
(pinned at v1.7.19) and an ANTLR-derived JSON grammar: a sanitizer-
instrumented harness classifies every generated input as accepted, rejected,
or crashing; a campaign layer reduces a run of the harness to three
observable, coverage-free proxies -- acceptance rate, the number of
distinct structural shapes produced, and the number of distinct rejection
reasons seen -- and a refinement loop feeds those proxies, together with the
grammar text, back to an LLM each iteration and asks it to propose a revised
Hypothesis strategy. Because the proposer is an LLM and its output is
executable code, every proposal is statically validated -- AST-checked for
disallowed imports and calls, and sanity-checked that it actually produces
well-formed output -- before it is ever run against the target binary.

We evaluate this pipeline with a seeded, five-run-per-arm comparison against
a static baseline generator, and separately by reusing the entire pipeline,
unmodified, against a second, independent parser (parson) to test whether
the approach generalizes beyond one target and one set of grammar
adaptations. Concretely, this paper contributes:

\begin{itemize}
  \item A black-box, coverage-free feedback design -- acceptance rate,
    structural diversity, and rejection signatures -- shown sufficient to
    steer LLM-driven refinement of a grammar-based generator without any
    coverage instrumentation.
  \item A pipeline that treats LLM-authored code as untrusted by
    construction: every proposed generator is statically sandboxed and
    behaviorally sanity-checked before it ever executes against a real
    target, independent of which LLM produced it.
  \item A seeded, reproducible, five-run-per-arm empirical comparison
    showing that refinement raises mean acceptance rate from 57.6\% to
    96.8\% relative to a static baseline on the same target. Structural
    fingerprints and rejection signatures are reported as descriptive
    signals, with explicit caveats that unequal per-arm budgets prevent
    direct comparison of their aggregate counts.
  \item A second, independently run target (parson), reusing every pipeline
    component unchanged, including five real refinement iterations and an
    honestly argued negative result -- no crash found across 3{,}000
    sanitizer-instrumented executions -- supported by manual review of the
    highest-risk code paths rather than left as an unexplained absence.
  \item An empirical account of grammar-implementation mismatches across
    both targets under one shared formal grammar -- for example, duplicate
    JSON object keys accepted by cJSON and rejected outright by parson --
    which we argue are themselves a source of fuzzing-relevant signal that
    a purely formal grammar specification cannot supply on its own.
  \item Reproducibility tooling for the generated-input stream specifically
    (a seeding shim over Hypothesis's internal random state, per-run
    manifests, and hash-identified harness binaries), together with an
    explicit statement of what that reproducibility does and does not cover.
\end{itemize}

The rest of the paper is organized as follows. Section~\ref{sec:related-work}
situates this work relative to coverage-guided, grammar-based, and
LLM-assisted fuzzing. Section~\ref{sec:methodology} describes the pipeline,
the proxy signal, and the sandboxing of LLM-proposed generators in detail.
Section~\ref{sec:evaluation} reports the cJSON comparison and the parson
bonus target. Section~\ref{sec:discussion} discusses limitations, threats to
validity, and what the negative result on parson does and does not show.
Section~\ref{sec:conclusion} concludes.
